\section{Star transitions and two-centre comparisons}\label{sec:blowup-comparisons}

We give the complete calculations for Theorem~\ref{branch:stars} and
Proposition~\ref{treemix:two-centres}, retaining the notation of
Sections~\ref{sec:mixed-branching} and~\ref{sec:mixed-tree-threshold}.

\subsection{Proof of the star transition}
\begin{proof}
The geometric assertions follow from Proposition~\ref{mixed:geometry}.
The unassigned divisor degree at every leaf is strictly less than one
by the local identity in Lemma~\ref{mixed:flux}. The tree criterion,
Theorem~\ref{mixed:tree}, therefore reduces stability to $A_m<1$.
This reduction includes $m=1$, when both vertices are leaves.

Use affine coordinates $a_i=1+\langle x,e_i\rangle$ and coordinates
$y_i=\langle x,u_i\rangle$. The anticanonical polytope is
\[
 a_i\ge0,\qquad \sum_{i=1}^m a_i\le m+1,\qquad
 -\min(a_i,1)\le y_i\le1.
\]
Define, for $a\ge0$ and $t>0$,
\begin{equation}\label{branch:star-volume}
 \ell(a)=1+\min(a,1),\qquad
 Z_m(t)=
 \int_{\substack{a_i\ge0\\\sum_i a_i\le t}}
       \prod_{i=1}^m\ell(a_i)\,da_1\cdots da_m.
\end{equation}
The lattice volume of the polytope is $Z_m(m+1)$. Eliminating one
coordinate from $\sum_i a_i=t$ gives the lattice measure on its central
facet, whose volume is $Z'_m(t)$. Consequently
\begin{equation}\label{branch:central-degree}
 A_m=\frac{Z'_m(m+1)}{Z_m(m+1)}.
\end{equation}
There is no Euclidean normal-length factor in this formula.

For $t>m$, the integral has the finite expression
\begin{equation}\label{branch:finite-sum}
 Z_m(t)=
 \sum_{j=0}^{m}(-1)^j\binom mj
 \sum_{h=0}^{m-j}\binom{m-j}{h}
 \frac{(t-j)^{2m-h}}{(2m-h)!}.
\end{equation}
Indeed, the Laplace transform of $\ell$ is
$(\lambda+1-e^{-\lambda})/\lambda^2$. The simplex integral in
\eqref{branch:star-volume} has transform
$(\lambda+1-e^{-\lambda})^m/\lambda^{2m+1}$; expanding the numerator
and inverting each term proves \eqref{branch:finite-sum}.
Differentiating that finite sum gives the exact values of $1-A_m$ in
Table~\ref{branch:finite-signs}. They prove stability for $1\le m\le5$
and $A_m>1$ for $6\le m\le14$. The entries are comparisons of the
individual divisor degree $A_m$ with one. In particular, the entry at
$m=1$ is not the weight deficit of the line $\CC e_0$, which then also
contains $e_1$.

\begin{table}[ht]
\centering
\caption{Exact values of $1-A_m$, obtained from
\eqref{branch:finite-sum} and its derivative at $t=m+1$.}
\label{branch:finite-signs}
\begin{adjustbox}{max width=\textwidth,center}
\begin{tabular}{r@{\qquad}l}
\toprule
$m$ & $1-A_m$\\
\midrule
\csname @@input\endcsname tables/mixed_star_signs.tex
\bottomrule
\end{tabular}
\end{adjustbox}
\end{table}

It remains to prove $A_m>1$ for all $m\ge15$. Set $t=m+1$ and
\[
 J_m(t)=
 \int_{\substack{a_i\ge0,\ \sum_i a_i\le t\\0\le a_1\le1}}
 a_1\prod_{i=2}^m\ell(a_i)\,da_1\cdots da_m.
\]
Rescale the simplex in \eqref{branch:star-volume} to a fixed simplex
and differentiate. Since $\ell'(a)=1$ on $(0,1)$ and is zero on
$(1,\infty)$, symmetry among the coordinates gives
\begin{equation}\label{branch:star-euler}
 tZ'_m(t)=mZ_m(t)+mJ_m(t).
\end{equation}
The kink of $\ell$ occurs on a set of measure zero and does not
contribute an additional term.

For every integer $p\ge1$, $s>0$, and $0<\theta\le1$, the inequalities
$\theta\ell(a)\le\ell(\theta a)\le\ell(a)$ imply
\[
 \theta^{2p}Z_p(s)\le Z_p(\theta s)\le\theta^p Z_p(s).
\]
Using the upper bound for $Z_m$ and the lower bound for $J_m$ yields
\begin{align*}
 Z_m(t)
 &=\int_0^t\ell(a)Z_{m-1}(t-a)\,da\\
 &\le2Z_{m-1}(t)\int_0^t(1-a/t)^{m-1}\,da
   =\frac{2t}{m}Z_{m-1}(t),\\
 J_m(t)
 &=\int_0^1aZ_{m-1}(t-a)\,da\\
 &\ge Z_{m-1}(t)\int_0^1a(1-a/t)^{2(m-1)}\,da.
\end{align*}
For $0\le a\le1$, the inequality
$\log(1-x)\ge-x/(1-x)$ gives
\[
 2(m-1)\log(1-a/t)
 \ge-\frac{2(m-1)a}{t-a}\ge-2a.
\]
Hence
\[
 \int_0^1a(1-a/t)^{2(m-1)}\,da
 \ge\int_0^1ae^{-2a}\,da
 =\frac{1-3e^{-2}}4>\frac17.
\]
The last strict inequality follows from $e^2>7$: the terms of its
power series through degree four sum to seven and the remaining terms
are positive. It follows that $J_m(t)/Z_m(t)>m/(14t)$. Substituting in
\eqref{branch:star-euler}, with $t=m+1$, proves
\begin{equation}\label{branch:star-tail}
 A_m-1
 =\frac{mJ_m(t)/Z_m(t)-1}{m+1}
 >\frac{m^2-14m-14}{14(m+1)^2}>0
 \qquad(m\ge15).
\end{equation}
The numerator on the right equals one at $m=15$ and is increasing
thereafter.

For $m\ge6$, the ray $e_0$ is the only ray on its line. The other
central rays are not collinear with the negative sum of at least two
basis vectors, and every other ray has a nonzero leaf component.
Thus $\CC e_0$ has normalized ray weight $A_m>1$ and gives the stated
strict destabilizer. Together with the finite table and the tree
criterion, this proves the theorem.
\end{proof}

\begin{corollary}\label{branch:hn}
For $m\ge6$, the first term of the Harder--Narasimhan filtration of
$T_{Y_m}$ is the line bundle $\mathcal O_{Y_m}(D_{e_0})$.
\end{corollary}

\begin{proof}
Put $A=A_m$, and let $B,C$ be the normalized degrees of $-u_i,s_i$,
respectively. The degrees of $u_i,e_i$ are $B-C,A-C$, so
$(m+1)A+2mB-mC=2m$. On the simplex
$\Delta=\{a_i\ge0:\sum_i a_i\le m+1\}$, put
$F(a)=\prod_{j\ne i}\ell(a_j)$ and $Z=Z_m(m+1)$. The leaf facets give
\[
 B=\frac{\int_\Delta F\,da}{Z},\qquad
 C=\frac{\int_{\Delta\cap\{a_i\le1\}}F\,da}{Z}.
\]
Since $\ell\le2$, strictly on an open set, $B>1/2$.
Also $Z(1-2B+C)=\int_{\Delta\cap\{a_i\le1\}}a_iF\,da>0$.
Thus $2B-C<1$, and positivity of $B-C$ gives $C<2B-C<1$.

The first Harder--Narasimhan term is unique, hence preserved by the torus
and by arm permutations. Its generic subspace is ray-spanned: otherwise
taking the span of its contained rays preserves its positive degree
and decreases its rank, contradicting maximality of slope.
Define
\[
 L=\CC e_0,\quad
 \mathcal E=\operatorname{span}(e_i),\quad
 \mathcal U=\operatorname{span}(u_i),\quad
 \mathcal S=\operatorname{span}(s_i).
\]
An invariant ray-spanned subspace is generated by complete ray orbits.
The two leaf-sign orbits have the same span; any two of
$\mathcal E,\mathcal U,\mathcal S$ span the ambient space, and
$L\subset\mathcal E$. Thus the nonzero proper possibilities are
$L,\mathcal E,\mathcal U,\mathcal S,L+\mathcal U,L+\mathcal S$.
Their normalized slopes at the three rank-$m$ spaces are
\[
 \widehat\mu(\mathcal E)=2-2B<1,\qquad
 \widehat\mu(\mathcal U)=2B-C<1,\qquad
 \widehat\mu(\mathcal S)=C<1.
\]
The last two spaces have slopes
$(A+m(2B-C))/(m+1)$ and $(A+mC)/(m+1)$, both strictly below $A>1$.
Hence this first term has generic subspace $L$. It is saturated and
rank one, hence reflexive and a line bundle on the smooth variety.
The induced filtration jumps only at $e_0$, with multiplicity one,
so this line bundle is $\mathcal O_{Y_m}(D_{e_0})$.
\end{proof}

Rank-one maximal destabilizing subsheaves also occur in the
four-dimensional examples of Reynolds~\cite[Appendix~A, p.~68]{Reynolds2023}.


\subsection{Proof for two branching vertices}
\begin{proof}
For the unit leaf input $\ell$, the incoming function supplied by a
degree-$d$ branching vertex with its $d-1$ pendant leaves is
\[
 g_d(a)=[z^d]e^{(1+a)z}Q_\ell(z)^{d-1}\qquad(0\le a\le1),
\]
by Lemma~\ref{treemix:coefficients}. It is constant above one. The three
polynomials are
\begin{align*}
 g_3(a)&=\frac{137+155a+60a^2+8a^3}{12},\\
 g_4(a)&=\frac{5217+5965a+2610a^2+520a^3+40a^4}{120},\\
 g_5(a)&=\frac{122889+141357a+66120a^2+15760a^3+1920a^4+96a^5}{720}.
\end{align*}
The degree-$d$ centre receives $T^{L-1}g_e$ from the connecting path.
For any positive input $p$, write $c=p(1)$,
$i=\int_0^1p(a)\,da$, and $j=\int_0^1ap(a)\,da$. Then
\begin{equation}\label{treemix:transfer-coefficients}
 T(p)(a)=(2i-j+c/2)+(i+c)a+(c/2)a^2.
\end{equation}
All coefficients are positive, since $2i-j>0$.

First consider the degree-two vertices of the connecting path. The
polynomials $g_d$ are convex on $[0,1]$, and their endpoint differences are
\[
 g_3(1)-g'_3(1)=\frac{61}{12},\qquad
 g_4(1)-g'_4(1)=\frac{1447}{120},\qquad
 g_5(1)-g'_5(1)=\frac{3821}{144}.
\]
The tangent-line inequality at one gives $g_d(a)\ge a g_d(1)$.
Equation~\eqref{treemix:transfer-coefficients} also gives
\[
 T(p)(a)-aT(p)(1)=(1-a)(2i-j+c/2-ac/2)\ge0.
\]
Every input at a connecting degree-two vertex therefore satisfies
\eqref{mixed:incoming-bound}. The two-input calculation in the proof
of Theorem~\ref{mixed:paths} again gives $A_v\le24/25<1$ at each of
these vertices, regardless of $L$.

It remains to treat the two branching vertices. For $d\in\{3,4,5\}$
define the linear functional
\begin{equation}\label{treemix:D}
 D_d(p)=[z^d]e^zQ_\ell(z)^{d-1}Q_p(z)
       -[z^{d-1}]e^zQ_\ell(z)^{d-1}Q_p(z).
\end{equation}
It equals $V-F$ at the degree-$d$ centre, so positivity is equivalent to
$A_v<1$. When a monomial $a^k$ is used as an input, it is understood to
be continued constantly with value one above one. Then
\[
 T(a^k)=\left(\frac2{k+1}-\frac1{k+2}+\frac12\right)
        +\left(\frac1{k+1}+1\right)a+\frac12a^2.
\]
The required values of $D_d(T(a^k))$ are printed in
Table~\ref{treemix:transfer-signs}. They are exact evaluations of
\eqref{treemix:Q} and \eqref{treemix:D}.

\begin{table}[ht]
\centering
\caption{Positive values of $D_d(T(a^k))$ for the capped monomial inputs.}
\label{treemix:transfer-signs}
\begin{adjustbox}{max width=\textwidth,center}
\begin{tabular}{rrrrrrr}
\toprule
$d$ & $k=0$ & $k=1$ & $k=2$ & $k=3$ & $k=4$ & $k=5$\\
\midrule
3 & $4477/360$ & $611/80$ & $1483/240$ & $7909/1440$ & $367/72$ & $24403/5040$\\
4 & $362/15$ & $19847/1440$ & $15361/1440$ & $44209/4800$ & $15077/1800$ & $79019/10080$\\
5 & $102241/3360$ & $725819/60480$ & $398891/60480$ & $55289/13440$ & $410239/151200$ & $15739/8640$\\
\bottomrule
\end{tabular}
\end{adjustbox}
\end{table}

For $L=2$, the input before the final application of $T$ is $g_e$, of
degree at most five with positive coefficients. For $L\ge3$, it is
$T^{L-2}g_e$, a quadratic with positive coefficients by
\eqref{treemix:transfer-coefficients}. In either case it is a nonnegative linear combination of the capped
monomials $1,a,\ldots,a^5$, with at least one positive coefficient.
Linearity and Table~\ref{treemix:transfer-signs} prove
$D_d(T^{L-1}g_e)>0$ for every $L\ge2$. Interchanging $d,e$ proves the
inequality at the other centre.

For $L=1$, the inputs are $g_e,g_d$ directly. The six possibilities up
to interchange are listed in Table~\ref{treemix:adjacent-signs}.
Their strictly positive values follow from the displayed polynomials
and coefficient formula \eqref{treemix:D}, divided by the corresponding
positive volumes.
\begin{table}[ht]
\centering
\caption{The smaller normalized deficit at the two adjacent branching vertices.}
\label{treemix:adjacent-signs}
\begin{adjustbox}{max width=\textwidth,center}
\begin{tabular}{cl}
\toprule
$(d,e)$ & $\min(1-A_{\mathrm{left}},1-A_{\mathrm{right}})$\\
\midrule
$(3,3)$ & $104355/994027$\\
$(3,4)$ & $2643419/54754807$\\
$(3,5)$ & $14018811/1317077851$\\
$(4,4)$ & $97814081/2152964421$\\
$(4,5)$ & $1281377927/155298531575$\\
$(5,5)$ & $8349076421/1244343460105$\\
\bottomrule
\end{tabular}
\end{adjustbox}
\end{table}

Every pendant leaf satisfies $A_v<1$ by Lemma~\ref{mixed:flux}.
Theorem~\ref{mixed:tree} now proves stability. The remaining assertions
follow from Proposition~\ref{mixed:geometry}.
\end{proof}
